% relatedwork.tex -- Related Work

\paragraph{Protein sequence generation.}
Profile hidden Markov models (pHMMs)~\cite{hmmer3,kroghHiddenMarkovModels1994} have long served as the standard tool for protein family modeling, capturing positional emission and transition probabilities from MSAs. While pHMMs can emit sequences, they model only single-site statistics and miss the pairwise and higher-order couplings critical for maintaining tertiary structure. Direct coupling analysis (DCA) and Potts models~\cite{monacoProteinsMutualInformation2014,ekeberg2013improved} capture pairwise epistasis but are generative only in the Boltzmann sense, requiring expensive MCMC sampling with learned couplings.

Deep generative models have transformed protein design. DeepSequence~\cite{riesselman2018deep} uses a VAE to learn a latent representation of protein families for fitness prediction and generation. ProtGPT2~\cite{ferruzProtGPT2DeepUnsupervised2022} fine-tunes a GPT-2 language model on UniRef50, generating full-length sequences autoregressively. ProGen~\cite{madaniProGenLanguageModeling2023} conditions on taxonomy and function tags for controllable generation. MSA Transformer~\cite{raoMSATransformer2021} processes entire alignments via tied row attention, learning inter-sequence dependencies. EvoDiff~\cite{alamdariProteinGeneration2023} applies discrete diffusion to protein sequences, achieving unconditional and MSA-conditioned generation. ESM-2~\cite{linEvolutionaryScaleModeling2023} and ESMFold provide large-scale protein language models whose representations enable structure prediction and design.

All of these methods require substantial training data---typically thousands to millions of sequences---and learned parameters that must be optimized. When the family of interest has $K<200$ members, the ratio of learnable parameters to training examples becomes unfavorable, leading to overfitting or mode collapse. Our approach requires no training: the memory matrix is the MSA itself, and the score function is analytic.

\paragraph{Modern Hopfield networks and attention.}
Ramsauer et al.~\cite{ramsauerHopfieldNetworksAll2021} showed that transformer-style attention is equivalent to one step of the retrieval dynamics on the modern Hopfield energy with a log-sum-exp interaction function. This connection, building on the dense associative memory framework of Krotov and Hopfield~\cite{krotovDenseAssociativeMemory2016}, provides exponential storage capacity and continuous-state retrieval. The Energy Transformer~\cite{hooverEnergyTransformer2024} exploits this energy for discriminative tasks, but does not sample from the associated Boltzmann distribution. Our prior work~\cite{alswaidanVarnerStochasticAttention2025} introduced stochastic attention and validated it on synthetic data, MNIST digits, financial time series, and a single protein family (RRM). The present work extends this framework to the protein domain with multiple families, structure-level validation, scaling analysis, and protein-specific baselines.

\paragraph{Langevin dynamics and score-based generation.}
Score-based generative models~\cite{songGenerativeModelingEstimating2019,songScoreBasedGenerativeModeling2021} and denoising diffusion models~\cite{hoDenoisingDiffusionProbabilistic2020} generate samples by learning a neural score function $\nabla \log p(\mathbf{x})$ and running Langevin or reverse-SDE dynamics. These methods achieve state-of-the-art sample quality but require a trained score network. In stochastic attention, the score is exact: $\nabla \log p_\beta(\boldsymbol{\xi}) = \beta(\mathbf{T}(\boldsymbol{\xi}) - \boldsymbol{\xi})$, where $\mathbf{T}$ is the softmax attention map. This eliminates score-estimation error and permits convergence guarantees inherited from the ULA literature~\cite{durmusNonasymptoticAnalysisUnadjusted2017,dalalyanTheoreticalGuaranteesApproximate2017}.
